\section{Lineage and Context}
\label{sec:doctrine-lineage}

\subsection{2016 Language-Model Lesson}
\label{sec:doctrine-lineage-lm}

The thesis lineage originates in a 2016 language-model experiment whose finding is stated
concisely in the lineage record: local text imitation does not conserve consequence. A system
that learns to predict the next token in a sequence of human activity records learns the
surface form of consequential behavior without acquiring any mechanism to validate, refuse, or
receipt-bear that behavior. Prediction is not coordination. Pattern-matching over logs
produces outputs that resemble lawful process behavior without being lawful process behavior.

This finding is the epistemic seed of the Blue River Dam doctrine. If a system cannot
distinguish between a lawful process trace and a statistically plausible imitation of one,
then no downstream consumer of that system's outputs can rely on them as process evidence.
The admission layer --- the typed, refusal-capable, receipt-bearing upstream gate --- is the
architectural response to this insufficiency. It provides what text imitation cannot: a named
witness, a specific violated law on refusal, and a replayable receipt chain proving that the
governing act occurred upstream of the data.

The 2016 lesson is not anti-language-model in character. It is a precision claim: language
models are valid downstream consumers of process truth. They may explain, route, compare,
summarize, and assist. What they may not do is adjudicate process truth. The Blue River Dam
doctrine, as stated in \texttt{BLUE\_RIVER\_DAM.md}, draws this line explicitly: ``Agents
may NOT adjudicate process truth. The kernel adjudicates; the agent cites deterministic
artifacts. No artifact, no claim.''

\subsection{Enterprise Process Gap}
\label{sec:doctrine-lineage-enterprise}

The enterprise process gap is the structural consequence of the 2016 lesson operating at
organizational scale. Enterprises accumulate activity records in logs, dashboards, audit
trails, and BI platforms that cannot prove a lawful process happened. They issue board-level
claims, M\&A diligence representations, and regulatory submissions built on data that has no
admission gate, no named witness, no typed refusal surface, and no replayable receipt.

The doctrine corpus characterizes this gap through the lens of the five maturity levels
defined in \texttt{PROCESS\_INTELLIGENCE\_DEFINED.md} and \texttt{BLUE\_RIVER\_DAM.md}. A
Level 1 organization records activity; it has no process mining capability, no formal process
model, no conformance checking, and no refusal capability. The vast majority of enterprise
software operates at Level 1 or near it. The sediment accumulates; the governing acts go
unwitnessed.

The enterprise process gap is not closed by adding a monitoring dashboard or an observability
layer. These are downstream tools. The gap is closed only by moving the admission and refusal
gate upstream --- before execution, not after it. This is the Blue River Dam position. The
dam must be upstream. Everything downstream of the dam is sediment.

\subsection{Relationship to the Chatman Equation}
\label{sec:doctrine-lineage-chatman-equation}

The Chatman Equation, expressed in the SPR thesis file as $R \vdash A = \mu(O^*)$, and in
the Blue River Dam operating equation as $\kappa(\rho(\alpha(\mu(O^*)))) \to \text{ALIVE}
\mid \text{PARTIAL} \mid \text{REFUSED}$, is the formal generalization of the 2016 lesson.
It states that admissible process authority $A$ is the result of manufacturing ($\mu$) from
the knowledge corpus $O^*$, under the refusal gate $R$.

The four-operator chain $\kappa \circ \rho \circ \alpha \circ \mu$ composes the
manufacturing stack:
\begin{itemize}
  \item $\mu(O^*)$ --- manufacture from the knowledge corpus: papers, standards, type laws,
    and doctrine files are the raw material.
  \item $\alpha(\cdot)$ --- actuate: the manufactured knowledge must be made consequential via
    the MAPE-K autonomic loop, not merely stored.
  \item $\rho(\cdot)$ --- emit evidence: actuation must produce typed, witnessed, receipted
    evidence artifacts, not narrative outputs.
  \item $\kappa(\cdot)$ --- gate the outcome: every result is classified as ALIVE (all gates
    met), PARTIAL (some gates met, gaps documented), or REFUSED (gate violation, named law
    cited).
\end{itemize}

No silent success is permitted under this equation. A result that does not achieve ALIVE must
be classified PARTIAL or REFUSED with explicit gap documentation. The sealed
\texttt{PROCESS\_INTELLIGENCE\_ALIVE\_001} checkpoint and the honest reporting of 15
algorithm families with execution missing are both direct applications of this equation.

\subsection{Relationship to Receipts and Replay}
\label{sec:doctrine-lineage-receipts}

The receipt and replay disciplines are the operationalization of the Chatman Equation's
evidence-emission requirement. A receipt, as defined in \texttt{RECEIPT\_DOCTRINE.md}, is
not a hash. A hash proves integrity of bytes. A receipt proves that a named law was applied,
by a named witness, to a typed result, within a lawful object lifecycle.

\begin{quote}
$\text{Receipt}\langle T, W\rangle$: $T$ is the result type; $W$ is the witness marker ---
a zero-sized type naming a specific paper, standard, or law (e.g., \texttt{Ocel20},
\texttt{Xes1849}, \texttt{WfNetSoundnessPaper}).
\end{quote}

Two receipts with different witnesses are different types at compile time. They cannot be
confused or substituted. This is the type-theoretic translation of the audit requirement:
a claim that something was done under OCEL 2.0 witness cannot be satisfied by a receipt
issued under XES 1849-2016 witness.

The replay discipline extends this: a receipt that cannot be replayed is narration. The
BLAKE3 receipt chaining formula $\text{Receipt}_n = \text{BLAKE3}(\text{Receipt}_{n-1}
\mid\mid \text{new\_action} \mid\mid \text{new\_state} \mid\mid \text{signature})$ produces
a chain of receipts that can be independently re-derived by any auditor possessing the
original inputs. This is the machine-verifiable audit trail that the enterprise process gap
currently lacks.

The lineage from receipts to the broader thesis is direct: receipts and replay are the
mechanism by which the ggen manufacturing machinery, the wasm4pm execution authority, and the
downstream capital-flow and settlement surfaces all connect to an auditable chain of process
truth. Without receipts, every downstream claim is an assertion. With receipts, every
downstream claim is a proof.
